 %-------------------------------------- COMIENZA descripción del caso de uso.

  \begin{UseCase}{CU04}{Consultar Detalle de Mascota}{
      Permite al Propietario visualizar toda la información detallada de una de sus mascotas, incluyendo su perfil básico, historial médico (vacunas,
  enfermedades, alergias), desparasitaciones y documentos asociados.
  }
      \UCitem{Versión}{\color{Gray}1.0}
      \UCitem{Autor}{\color{Gray}Ugalde Tellez Aaron}
      \UCitem{Supervisa}{\color{Gray}Vazquez Villagran Jorge}
      \UCitem{Actor}{\hyperlink{Propietario}{Propietario}}
      \UCitem{Propósito}{Consultar de forma centralizada y completa toda la información relevante de una mascota para su correcta gestión, cuidado y para
  compartirla con el personal del hotel.}
      \UCitem{Entradas}{
          \begin{itemize}
              \item ID de la mascota a consultar (implícito en la navegación).
              \item Credenciales de autenticación del Propietario (token de sesión).
          \end{itemize}
      }
      \UCitem{Origen}{Mouse / Pantalla táctil}
      \UCitem{Salidas}{
          \begin{itemize}
              \item Pantalla con la información detallada de la mascota.
              \item Lista de vacunas, enfermedades, alergias, desparasitaciones y documentos.
          \end{itemize}
      }
      \UCitem{Destino}{Pantalla}
      \UCitem{Precondiciones}{
          \begin{itemize}
              \item El Propietario ha iniciado sesión en el sistema.
              \item El Propietario se encuentra en la \IUref{IU-PetList}{Lista de Mascotas}.
          \end{itemize}
      }
      \UCitem{Postcondiciones}{Ninguna. Este caso de uso es de sólo consulta y no modifica el estado del sistema.}
      \UCitem{Errores}{
          \begin{itemize}
              \item \textbf{E1}: La mascota no existe.
              \item \textbf{E2}: El usuario no está autorizado para ver la mascota.
              \item \textbf{E3}: Error de conexión o fallo del servidor.
          \end{itemize}
      }
      \UCitem{Tipo}{Primario}
  \end{UseCase}
  %--------------------------------------
  % Trayectoria Principal
  %--------------------------------------
  \begin{UCtrayectoria}
      \UCpaso[\UCactor] Selecciona una mascota específica desde la \IUref{IU-PetList}{Lista de Mascotas}. \label{CU04-Start}
      \UCpaso El sistema navega a la \IUref{IU-PetDetail}{Pantalla de Detalle de Mascota}.
      \UCpaso Solicita al servidor toda la información asociada a la mascota seleccionada (GET /api/pets/:id, GET /api/pet-vaccinations/:id, etc.).
      \UCpaso El sistema verifica que la sesión del usuario sea válida con base en la regla \BRref{BR-AuthCheck}{Verificación de Autenticación}. \Trayref{A}.
      \UCpaso El sistema verifica que la mascota solicitada pertenezca al propietario autenticado con base en la regla \BRref{BR-PetOwnership}{Verificación de Propiedad de Mascota}. \Trayref{B}.
      \UCpaso El sistema obtiene de la base de datos el perfil de la mascota y su historial médico relacionado (vacunas, enfermedades, etc.). \Trayref{C}.
      \UCpaso Muestra la información completa y organizada en secciones dentro de la \IUref{IU-PetDetail}{Pantalla de Detalle de Mascota}.
  \end{UCtrayectoria}
  
  %--------------------------------------
  % Trayectorias Alternativas
  %--------------------------------------
  
  \begin{UCtrayectoriaA}{A}{Sesión inválida o expirada}
      \UCpaso[CU04-\ref{CU04-Start}] El sistema detecta que el token de autenticación es inválido o no existe.
      \UCpaso Muestra el mensaje de error {\bf MSG-Error} ``Su sesión ha expirado. Por favor, inicie sesión nuevamente.''.
      \UCpaso[] Redirige al usuario a la \IUref{IU-Login}{Pantalla de Inicio de Sesión}.
      \UCpaso[] Termina el caso de uso.
  \end{UCtrayectoriaA}
  
  \begin{UCtrayectoriaA}{B}{Mascota no encontrada o no autorizada}
      \UCpaso[CU04-\ref{CU04-Start}] El sistema determina que la mascota no existe (Error 404) o no pertenece al propietario (Error 403).
      \UCpaso Muestra el mensaje de error {\bf MSG-Error} ``Mascota no encontrada''.
      \UCpaso[] Redirige al usuario a la \IUref{IU-PetList}{Lista de Mascotas}.
      \UCpaso[] Termina el caso de uso.
  \end{UCtrayectoriaA}
  
  \begin{UCtrayectoriaA}{C}{Fallo inesperado del sistema}
      \UCpaso Ocurre un error interno en el servidor al intentar consultar la base de datos (Error 500).
      \UCpaso Muestra el mensaje de error {\bf MSG-Error} ``Error al cargar los detalles de la mascota''.
      \UCpaso[\UCactor] Oprime el botón \IUbutton{Aceptar} o recarga la página.
      \UCpaso[] El sistema permanece en una vista de error, permitiendo al usuario reintentar la navegación.
  \end{UCtrayectoriaA}
  
  %--------------------------------------
  % Puntos de extensión
  %--------------------------------------
  \subsection{Puntos de extensión}
  \UCExtenssionPoint{
      % Cuando:
      El Propietario desea modificar la información básica de la mascota.
  }{
      % Durante la región:
      Paso 7 (Visualización de los detalles de la mascota).
  }{
      % Casos de uso a los que extiende:
      \hyperlink{CU-03}{CU-03 Ver todas mis mascotas}.
  }
  \UCExtenssionPoint{
      % Cuando:
      El Propietario desea registrar una nueva vacuna, enfermedad, alergia, desparasitación o documento.
  }{
      % Durante la región:
      Paso 7 (Visualización de los detalles de la mascota).
  }{
      % Casos de uso a los que extiende:
      \hyperlink{CU-05}{CU-05 Agregar Vacuna}, \hyperlink{CU-07}{CU-07 Agregar Alergia}, etc.
  }    %-------------------------------------- TERMINA descripción del caso de uso.